\chapter{仓室模型和神经网络模型的构建} \label{ch:models}


\section{具有年龄结构和疫苗接种的 SEIRV 仓室模型}

在 COVID-19 疫情暴发中，地区内的所有人分为五个仓室：易感者 $ S(t) $，潜伏者 $ E(t) $，感染者 $ I(t) $，康复者 $ R(t) $ 和接种疫苗者 $ V(t) $。
模型中的潜伏者包括被感染但由于病毒负载低而难以向其他个体传播的个体，而感染者指能向其他个体传播病毒的个体。
模型进一步假设，G1代表60岁以下的个体，G2表示60岁及以上的个体。显然，G2在 COVID-19 传播过程中面临着高风险，因为G2中患有高血压、糖尿病、高脂血症、心脏病、慢性炎症等基础疾病的人数较多。
因为不同年龄组和性别组的疫苗接种情况存在差异，所以模型还考虑了接种疫苗者的情况。

该年龄结构仓室模型如下所示：
\begin{equation}
	\label{eq:model}
	\begin{aligned}
		\dot{S}_1 & = \Lambda-(\beta_{11} I_1+\beta_{12} I_2) \dfrac{S_1}{N} -(\nu_1+\mu_1+g) S_1,            \\
		\dot{S}_2 & = g S_1-(\beta_{21} I_1+\beta_{22} I_2) \dfrac{S_2}{N}-\left(\nu_2+\mu_2\right) S_2,      \\
		\dot{E}_1 & = (\beta_{11} I_1+\beta_{12} I_2) \dfrac{S_1}{N}-\left(\alpha_1+\mu_1+g\right) E_1,     \\
		\dot{E}_2 & = g E_1+(\beta_{21} I_1+\beta_{22} I_2) \dfrac{S_2}{N}-\left(\alpha_2+\mu_2\right) E_2, \\
		\dot{I}_1 & = \alpha_1 E_1-\left(\gamma_1+d_1+\mu_1+g\right) I_1,                                  \\
		\dot{I}_2 & = gI_1+\alpha_2 E_2-\left(\gamma_2+d_2+\mu_2\right) I_2,                               \\
		\dot{R}_1 & = \gamma_1 I_1-\left(\mu_1+g\right) R_1,                                               \\
		\dot{R}_2 & = g R_1+\gamma_2 I_2-\mu_2 R_2,                                                        \\
		\dot{V}_1 & = \nu_1 S_1-\left(\mu_1+g\right) V_1,                                                  \\
		\dot{V}_2 & = g V_1+\nu_2 S_2-\mu_2 V_2 \text{.}
	\end{aligned}
\end{equation}

模型中的所有参数均假定为常数，参数的描述如表~\ref{tab:para-des}所示。

\begin{table}[H]
	\footnotesize
	\centering
	\begin{threeparttable}[c]
		\caption{模型的参数及其描述}
		\label{tab:para-des}
		\begin{tabular*}{1\textwidth}{@{\extracolsep{\fill}}cllrrl@{}}
			\toprule
			\multicolumn{1}{l}{参数} & \multicolumn{1}{l}{含义}       & \multicolumn{1}{l}{单位\tnote{f}}    & \multicolumn{1}{c}{值}                  & \multicolumn{1}{c}{范围}                     & 来源                                                 \\
			\midrule
			$\Delta T$                & 失觉期                          & d                               & 21                     & (17, 23)                  & ~\citep{Data_AD,SinaNews20210107,SinaNews20210111}                  \\
			$N$                       & 石家庄市总人口                      & 万                               & \num{1900}\tnote{c}    & \num{1900}                & ~\citep{Data_people}                                                \\
			$\Lambda$                 & 每日新增人口的比率                    & d$^{-1}$                        & \num{300.82}           & \num{1000}\tnote{d}       & ~\citep{Data_age_rate}                                              \\
			$g$                       & G1\tnote{a}到G2\tnote{b}的衰老速率 & d$^{-1}$                        & \num{1.36e-5}\tnote{e} & \num{4e-4}                & ~\citep{Data_g}                                                     \\
			$\beta_{11}$              & $ I_1 $引发的感染率\tnote{a}       & p$^{-1}$d$^{-1}$                & \num{0.1790}           & (0, 0.5)                  & ~\citep{Data_AD}                                                    \\
			$\beta_{12}$              & $ I_2 $引发的感染率\tnote{a}       & p$^{-1}$d$^{-1}$                & \num{0.1721}           & (0, 0.5)                  & ~\citep{Data_AD}                                                    \\
			$\beta_{21}$              & $ I_1 $引发的感染率\tnote{b}       & p$^{-1}$d$^{-1}$                & \num{0.2833}           & (0, 0.5)                  & ~\citep{Data_AD}                                                    \\
			$\beta_{22}$              & $ I_2 $引发的感染率\tnote{b}       & p$^{-1}$d$^{-1}$                & \num{0.2479}           & (0, 0.5)                  & ~\citep{Data_AD}                                                    \\
			$1 / \alpha_{1}$          & 平均潜伏期\tnote{a}               & d                               & 5                      & (\num{3.53},  \num{6.6})    & ~\citep{Data_incubation,Data_incubation_1,ZhuBMCMed2021}            \\
			$1 / \alpha_{2}$          & 平均潜伏期\tnote{b}               & d                               & 3                      & (\num{3.53},  \num{6.6})    & ~\citep{Data_incubation,Data_incubation_1,ZhuBMCMed2021}            \\
			$1 / \gamma_{1}$          & 平均康复期\tnote{a}               & d                               & 28                     & (\num{13.6},  \num{33.4})   & ~\citep{Data_gamma_0,Data_gamma_1}                                  \\
			$1 / \gamma_{2}$          & 平均康复期\tnote{b}               & d                               & 35                     & (\num{13.6},  \num{33.4})   & ~\citep{Data_gamma_0,Data_gamma_1}                                  \\
			$\nu_{1}$                 & 接种率\tnote{a}                 & d$^{-1}$                        & 0                      & 0                         & ~\citep{Data_v}                                                     \\
			$\nu_{2}$                 & 接种率\tnote{b}                 & d$^{-1}$                        & 0                      & 0                         & ~\citep{Data_v}                                                     \\
			$d_{1}$                   & SARS-CoV-2导致的死亡率\tnote{a}    & d$^{-1}$                        & \num{2.37e-7}          & (1, 3$) \times \num{e-7}  $ & ~\citep{Data_clinic,Data_death,Data_death1,Data_death2,Data_death3}                    \\
			$d_{2}$                   & SARS-CoV-2导致的死亡率\tnote{b}    & d$^{-1}$                        & \num{1.42e-7}          & (1, 4$) \times \num{e-7}  $ & ~\citep{Data_clinic,Data_death,Data_death1,Data_death2,Data_death3}                     \\
			$\mu_{1}$                 & 自然死亡率\tnote{a}               & d$^{-1}$                        & \num{2.37e-7}          & \num{2.37e-7}             & ~\citep{Data_mu}                                                    \\
			$\mu_{2}$                 & 自然死亡率\tnote{b}               & d$^{-1}$                        & \num{1.42e-7}          & \num{1.42e-7}             & ~\citep{Data_mu}                                                    \\
			$p$                       & 老龄化率\tnote{b}                & -                               & 19.38\%                & 19.38\%                   & ~\citep{Data_age_rate}                                              \\
			\bottomrule
		\end{tabular*}
		\begin{tablenotes}
			\footnotesize
			\item[a]{G1 表示60岁以下的年龄组；}
			\item[b]{G2 表示60岁及以上的年龄组；}
			\item[c]{石家庄市在COVID-19疫情期间的人口规模是 \num{1900} 万人；}
			\item[d]{每年新增 \num{1000} 万人除以 \num{365} 天等于每天新增 \num{300.82} 万人；}
			\item[e]{每年的衰老速率 \num{4e-4} 除以 \num{365} 天等于每天的衰老速率 \num{1.36e-5}；}
			\item[f]{d表示天，d$^{-1}$表示每天，p$^{-1}$d$^{-1}$表示每人每天。}
		\end{tablenotes}
	\end{threeparttable}
\end{table}


\section{LSTM 神经网络模型}

研究传染病流行的流行病学问题通常需要采用时间序列分析方法，这类方法基于对系统观察获得的数据进行建模和分析，以预测未来的传播趋势。
其中，一种强大的工具是深度学习中的长短时记忆网络（LSTM，Long Short-Term Memory），该网络已被广泛用于处理和预测各种时间序列问题，
包括流行病学数据的分析~\cite{Li2023AO,Xu2019RAA,Pillai2021SR}。

由 Hochreiter 和 Schmidhuber （1997）提出的 LSTM（Long Short-Term Memory）网络广泛用于解决具有长期相关特征的时间序列问题。
LSTM 网络是一种循环神经网络（Recurrent Neural Network，RNN）的变体，设计用来解决传统RNN存在的长期依赖问题。
LSTM引入了门控机制，能够更好地捕捉和管理序列数据中的长期依赖关系。

该模型的内部结构如图~\ref{fig:Inner_Structure}所示。

\begin{figure}[H]
	\centering
	\includegraphics[width=1\textwidth]{Inner_Structure}
	\caption{LSTM 模型的内部结构}
	\label{fig:Inner_Structure}
\end{figure}

该模型可用方程~\eqref{eq:LSTM}表示。

\begin{equation}
	\label{eq:LSTM}
	\left\{
	\begin{array}{l}
		f_t=\sigma\left(W_f \cdot\left[h_{t-1}, x_t\right]+b_f\right)             \\
		i_t=\sigma\left(W_i \cdot\left[h_{t-1}, x_t\right]+b_i\right)             \\
		\widetilde{C}_t=\tanh \left(W_C \cdot\left[h_{t-1}, x_t\right]+b_C\right) \\
		C_t=f_t * C_{t-1}+i_t * \widetilde{C}_t                                   \\
		o_t=\sigma\left(W_o \cdot\left[h_{t-1}, x_t\right]+b_o\right)             \\
		h_t=o_t * \tanh \left(C_t\right)
	\end{array}
	\right.
\end{equation}

方程~\eqref{eq:LSTM}包括以下主要组成部分：

1. 输入门（Input Gate）：

输入门决定哪些信息将被存储到细胞状态（cell state）中。它由以下几个步骤组成：

a. 计算遗忘门（Forget Gate）：通过 $f_t= \sigma(W_f \cdot [h_{t-1}, x_t] + b_f)$计算遗忘门的值，它决定了以前的细胞状态中哪些信息要被遗忘。
其中，$W_f$ 是遗忘门的权重矩阵，$h_{t-1}$ 是上一个时间步的隐藏状态，$x_t$ 是当前时间步的输入，$b_f$ 是偏置项，$\sigma$ 是sigmoid激活函数。

b. 计算输入门的值（Input Gate）： 输入门决定哪些新的信息将被添加到细胞状态中，即$i_t = \sigma(W_i \cdot [h_{t-1}, x_t] + b_i)$。

c. 计算候选细胞状态（Candidate Cell State）： 计算一个候选的新细胞状态 $\tilde{C}_t= \tanh(W_C \cdot [h_{t-1}, x_t] + b_C)$。

2. 更新细胞状态：使用遗忘门、输入门以及候选细胞状态来更新细胞状态：$C_t = f_t \cdot C_{t-1} + i_t \cdot \tilde{C}_t$。

3. 输出门（Output Gate）： 输出门决定哪些信息将被输出到隐藏状态以及最终的预测中。 计算输出门的值 $o_t = \sigma(W_o \cdot [h_{t-1}, x_t] + b_o)$。

4. 隐藏状态更新： 使用输出门和细胞状态来更新当前时间步的隐藏状态：$h_t = o_t \cdot \tanh(C_t)$。

这些方程描述了LSTM模型的基本运算过程。
LSTM通过门控机制，能够选择性地遗忘、存储和输出信息，从而更好地处理长序列依赖关系，适用于各种序列建模任务，如自然语言处理、时间序列预测等。
这些门控机制允许LSTM模型有效地学习和记忆序列中的重要信息。


本研究将基于莆田市的COVID-19感染数据，借助SEIRV传染病模型，通过调整传播概率、潜伏期、康复或死亡概率以及接触人数等参数，构建一个基本的训练数据集。
此后，利用LSTM时间序列模型，结合仓室模型的基础，来研究病毒传播趋势，并进行石家庄市COVID-19传播的预测。
这种结合深度学习和传染病建模的方法有望提供更准确的流行病学预测和决策支持。

